\section{Molecules and the lifted correction}\label{sec:molecules}

We now turn the selected tests into a Banach space. The norm must do
two jobs: control point differences so that the correction can be
lifted Lipschitzly, and remember lists of tests whose signed sums have
a common Lipschitz bound. The latter feature will rule out linear
\(c_0\).

\subsection{The quotient and its completion}

Let \(\mathcal M_0(B)\) be the real vector space of mass-zero finite
formal sums
\[
 m=\sum_{s\in B}a_s\delta_s,\qquad
 \#\{s:a_s\ne0\}<\infty,\qquad \sum_s a_s=0.
\]
For a real function \(u\) on \(B\), set
\(\langle u,m\rangle=\sum_s a_su(s)\). Constants pair to zero.
This follows the mass-zero molecule construction of Arens and Eells
\cite{AE56} and its later Lipschitz-free formulation \cite{GK03}.
The norm below is built from the selected classes
\eqref{eq:classes}; no universal property of the usual free space is
assumed.

For every triple \(d=(\mathcal L,C,p)\), where \(\mathcal L\) is one
of the listed classes, \(C\in\N\), and \(p=1+1/k\) for \(k\in\N\),
define
\begin{equation}
 \sigma_d(m)=
 \sup\left\{
 \left(\sum_{\nu=1}^l|\langle u_\nu,m\rangle|^p\right)^{1/p}:
 \begin{array}{l}
 l\geq0,\quad u_1,\ldots,u_l\in\mathcal L,\\
 \Lip\!\left(\sum_{\nu=1}^l e_\nu u_\nu\right)\leq C
 \ \text{for every }(e_\nu)\in\{-1,1\}^l
 \end{array}\right\}.                                      \label{eq:sigma}
\end{equation}
The empty list contributes zero. Lists may contain repetitions and
tests with the same representing label; their sole compatibility
condition is the displayed Lipschitz bound for every choice of signs.
Equivalently, a list is admissible precisely when
\begin{equation}
 \sum_{\nu=1}^l|u_\nu(s)-u_\nu(t)|\leq C\|s-t\|_D
                   \qquad(s,t\in B).
 \label{eq:admissible-increments}
\end{equation}
To obtain this inequality, choose signs matching the increments
\(u_\nu(s)-u_\nu(t)\); the converse follows from the triangle inequality.

Enumerate the triples by \(d=1,2,\ldots\), write \(C_d\) for their
integer component, and set
\begin{equation}
             \alpha_d=\frac{2^{-d}}{C_d^2},\qquad
 \|m\|_0=\left(\sum_{d\geq1}\alpha_d\sigma_d(m)^2\right)^{1/2}.
                                                               \label{eq:norm}
\end{equation}

The outer square sum and the exponents approaching one serve distinct
purposes in Section~\ref{sec:no-c0}. The square sum forces one fixed
seminorm to detect a subsequence of any hypothetical \(c_0\) basis.
For that fixed test class and Lipschitz bound, a second exponent
\(p'<p\) converts detection by a list into detection by one scalar
test. Both choices are part
of the obstruction, as well as of the norm construction.

\begin{lemma}[Exact completed norm]\label{lem:completion}
Formula \eqref{eq:norm} is a finite seminorm on \(\mathcal M_0(B)\).
Let \(N\) be its kernel, and let \(Z\) be the Banach completion of
\(\mathcal M_0(B)/N\) with the induced norm. Then \(Z\) is separable,
and
\begin{equation}
 \zeta:B\to Z,\qquad \zeta(s)=[\delta_s-\delta_0]             \label{eq:zeta}
\end{equation}
is \(1\)-Lipschitz with \(\zeta(0)=0\). Each \(\sigma_d\) descends
and extends continuously to a seminorm on \(Z\), and the exact formula
\begin{equation}
                    \|z\|_Z^2=\sum_{d\geq1}
                                \alpha_d\sigma_d(z)^2         \label{eq:completed-norm}
\end{equation}
holds for every \(z\in Z\).
\end{lemma}

\begin{proof}
For \(s,t\in B\) and an admissible list for \(d\), inequality
\eqref{eq:admissible-increments} bounds the sum of the absolute
increments. Since the \(\ell_p\) norm is no larger than this sum,
\begin{equation}
                 \sigma_d(\delta_s-\delta_t)
                         \leq C_d\|s-t\|_D.                  \label{eq:molecule-bound}
\end{equation}
For a fixed list, Minkowski's inequality makes its expression in
\eqref{eq:sigma} a seminorm in \(m\); the supremum of these expressions
is also a seminorm, once finite. Every \(m=\sum_s a_s\delta_s\)
of mass zero equals \(\sum_s a_s(\delta_s-\delta_0)\). Thus
\[
 \sigma_d(m)\leq C_d\sum_s|a_s|\|s\|_D.
\]
This proves finiteness both of \(\sigma_d\) and of \(\|m\|_0\),
because \(\sum_d\alpha_d C_d^2=\sum_d2^{-d}=1\).
Minkowski's inequality in \(\ell_2\), applied to the weighted
seminorms, proves subadditivity of \(\|\cdot\|_0\); homogeneity is
immediate. Its kernel is therefore a linear subspace, and quotienting
by it produces a norm.

Equation \eqref{eq:molecule-bound} and the same square sum give
\[
                   \|\delta_s-\delta_t\|_0\leq\|s-t\|_D.
\]
This proves the assertion about \(\zeta\). The ball \(B\) is separable
because \(D\) is a countable-coordinate \(c_0\) space. If \(B_0\) is
a countable dense subset, continuity of \(\zeta\) makes
\(\zeta(B_0)\) dense in \(\zeta(B)\). The real linear span of
\(\zeta(B)\) is the quotient molecule space; its rational span over
\(B_0\) is consequently dense in \(Z\). Hence \(Z\) is separable.

For molecules \(m,m'\),
\begin{equation}
 |\sigma_d(m)-\sigma_d(m')|
 \leq\sigma_d(m-m')\leq\alpha_d^{-1/2}\|m-m'\|_0.            \label{eq:sigma-continuity}
\end{equation}
In particular \(\sigma_d\) vanishes on \(N\), descends, and extends
continuously to \(Z\); subadditivity and homogeneity pass to limits.
Moreover the map
\[
 m\longmapsto(\sqrt{\alpha_d}\sigma_d(m))_{d\geq1}
\]
is \(1\)-Lipschitz into \(\ell_2\), since the coordinate differences
are bounded by \(\sqrt{\alpha_d}\sigma_d(m-m')\). If molecules
converge to \(z\in Z\), their images therefore converge in \(\ell_2\).
The \(d\)-th coordinate of the limit is
\(\sqrt{\alpha_d}\sigma_d(z)\) by continuity of each seminorm.
The norms of those images are exactly the molecule norms, which
converge to \(\|z\|_Z\). This proves \eqref{eq:completed-norm}.
\end{proof}

\subsection{Fixed-class duality}

The obstruction argument will modify a test through a pointwise
limit after a vector in \(Z\) has already been fixed. The next lemma
justifies that passage, and also retains the exact list formula after
completion.

\begin{lemma}[Dual tests and pointwise limits]\label{lem:duality}
Every test in \(\mathcal T_h\) or \(\mathcal V_h\) defines a bounded
linear functional on \(Z\) by its molecule pairing. For each fixed
class \(\mathcal L\), there is a finite common bound \(B_{\mathcal L}\)
for these dual norms. For each fixed triple \(d=(\mathcal L,C,p)\),
formula \eqref{eq:sigma} remains exact on \(Z\), with the extended
test pairings.

If \(u_n\in\mathcal L\) and \(u_n(s)\to u(s)\) for every \(s\in B\),
then the pointwise limit defines a bounded functional on \(Z\) of
norm at most \(B_{\mathcal L}\), and
\(\langle u_n,z\rangle\to\langle u,z\rangle\) for each fixed
\(z\in Z\). More precisely, for any molecule class \(m\) and any
\(z\in Z\),
\begin{equation}
 |\langle u_n-u,z\rangle|
 \leq|\langle u_n-u,m\rangle|
       +2B_{\mathcal L}\|z-m\|_Z.                            \label{eq:fixed-vector}
\end{equation}
\end{lemma}

\begin{proof}
Choose an integer \(C_{\mathcal L}\) no smaller than the uniform
Lipschitz bound in Lemma~\ref{lem:class-bounds}. Every singleton
\((u)\), \(u\in\mathcal L\), is admissible for the one triple
\(d_{\mathcal L}=(\mathcal L,C_{\mathcal L},2)\). Thus
\begin{equation}
 |\langle u,m\rangle|
 \leq\sigma_{d_{\mathcal L}}(m)
 \leq\alpha_{d_{\mathcal L}}^{-1/2}\|m\|_0.                 \label{eq:dual-bound}
\end{equation}
It vanishes on \(N\) and extends to \(Z^*\), uniformly with
\(B_{\mathcal L}=\alpha_{d_{\mathcal L}}^{-1/2}\).

For an admissible list \(\boldsymbol u=(u_1,\ldots,u_l)\) of a
fixed triple \(d\), its evaluation map
\(\Lambda_{\boldsymbol u}m=(\langle u_\nu,m\rangle)_{\nu=1}^l\)
into \(\ell_p^l\) has norm at most \(\alpha_d^{-1/2}\) on the
molecule quotient. It therefore extends with that bound to \(Z\).
The supremum of \(\|\Lambda_{\boldsymbol u}z\|_p\) over all admissible
lists is finite and \(\alpha_d^{-1/2}\)-Lipschitz as a seminorm in
\(z\): its change between \(z,z'\) is bounded by the same supremum
at \(z-z'\). This supremum agrees with \(\sigma_d\) on the dense
molecule quotient. It is therefore the continuous extension of
\(\sigma_d\), which proves the exact formula on \(Z\).

For a pointwise-convergent sequence, the pairings converge on every
finite molecule. Inequality \eqref{eq:dual-bound} passes to that
pointwise limit, so its molecule functional extends with norm at most
\(B_{\mathcal L}\). Applying the two common dual bounds to \(z-m\)
gives \eqref{eq:fixed-vector}. For a prescribed \(z\) and tolerance
\(\delta>0\), choose a molecule class \(m\) with
\(\|z-m\|_Z<\delta/(4B_{\mathcal L})\), and then choose \(n\) so large
that the finite pairing in \eqref{eq:fixed-vector} is less than
\(\delta/2\). This proves convergence on \(z\). The choice of \(m\)
and \(n\) may depend on \(z\); no convergence uniform over a moving
sequence of vectors is asserted.
\end{proof}

\begin{lemma}[The \(c_0\)-valued correction map]\label{lem:Q}
There is a bounded linear map \(Q:Z\to U\) such that
\begin{equation}
                         Q\zeta(s)=h(s)=(S_s-P)s
                         \qquad(s\in B).                    \label{eq:Qzeta}
\end{equation}
\end{lemma}

\begin{proof}
Define on formal molecules
\[
                  Q_0m=\sum_s a_s h(s)\in U.
\]
This finite sum lies in \(U\), not merely in \(\ell_\infty\), because
each \(h(s)\) lies in \(U\). Each coordinate function \(f_\gamma\)
belongs to \(\mathcal T_1\) through its zero-operation representation.
Choose an integer \(C_0\geq L_0\) and the triple
\(d_0=(\mathcal T_1,C_0,2)\). Every singleton \((f_\gamma)\) is
admissible, whence
\[
 \|Q_0m\|_U=\sup_\gamma|\langle f_\gamma,m\rangle|
 \leq\sigma_{d_0}(m)\leq\alpha_{d_0}^{-1/2}\|m\|_0.
\]
Thus \(Q_0\) vanishes on \(N\), descends, and extends to a bounded
linear map into the complete space \(U\). Since \(h(0)=0\), applying
it to \(\delta_s-\delta_0\) proves \eqref{eq:Qzeta}.
\end{proof}

\subsection{Radial extension}

The point map \(\zeta\) is defined on \(B\). Its radial extension
lifts the homogeneous shift correction on all of \(D\).

\begin{lemma}[Lipschitz lift]\label{lem:K}
Define
\begin{equation}
 K(0)=0,\qquad
 K(x)=\|x\|_D\,\zeta\!\left(\frac{x}{\|x\|_D}\right)
                         \quad(x\ne0).                       \label{eq:K}
\end{equation}
Then \(K:D\to Z\) is positively homogeneous and
\(\Lip(K)\leq3\). It satisfies the exact identity
\begin{equation}
                              QK=g-P.                        \label{eq:lift}
\end{equation}
\end{lemma}

\begin{proof}
Since \(\zeta(0)=0\) and \(\Lip(\zeta)\leq1\),
\(\|\zeta(s)\|_Z\leq1\) for \(s\in B\). Suppose
\(a=\|x\|_D\geq b=\|x'\|_D>0\), and write \(s=x/a\), \(s'=x'/b\).
Splitting at \(b\zeta(s)\) gives
\[
 \|K(x)-K(x')\|_Z
 \leq(a-b)\|\zeta(s)\|_Z+b\|\zeta(s)-\zeta(s')\|_Z
 \leq\|x-x'\|_D+b\|s-s'\|_D.
\]
The normalization estimate \eqref{eq:normalization} bounds the last term by
\(2\|x-x'\|_D\). If an endpoint is zero, use
\(\|K(x)\|_Z\leq\|x\|_D\). This proves the Lipschitz bound.
For \(x\ne0\), \eqref{eq:Qzeta}, \eqref{eq:h-definition}, and \eqref{eq:g-definition} give
\[
 QK(x)=\|x\|_D h(x/\|x\|_D)
      =S_{x/\|x\|_D}x-Px=g(x)-Px.
\]
All terms vanish at zero, and positive homogeneity follows directly
from \eqref{eq:K}.
\end{proof}

We have now obtained the exact data \(Q,K,g\) of
Lemma~\ref{lem:absorption}. No surjectivity of \(Q\) was used or
proved. The remaining task is to show that the completed space \(Z\)
has no linear copy of \(c_0\).
